% Status: unverified working notes; not included by manuscript.tex.
% Equations and prose require correction before integration into a chapter.
\section{Research Journey}

Instead of doing a typical literature review, I will explain how a computational material scientist ends doing research in operator theory.

So this journey starts with me deciding to teach at the University of the Philippines, and going from large R1 universties and national laboratories, I move to a place where I am teaching solid state physics.

So my having done a lot of math modelling and then implementing it in code for twenty years, you realize you don't really understand something if you cannot implement it into software.

So it start innocently enough,
$$
  \hat{H} = \hat{T} + \hat{V}
$$
where $\hat{H}$ is the quantum Hamiltonian operator, $hat{T}$ is the kinetic operator,and $\hat{V}=V(x)$.

Which we can turn into an eigenvalue/eigenvector problem by
\begin{equation}
  \langle{\mathbf{x}}|{\hat{H}}| \psi \rangle
  = E \langle x | \psi x \rangle
\end{equation}
which becomes
$$
  H \boldsymbol{\psi} = \varepsilon \boldsymbol{\psi}.
$$
